\documentclass[10pt,a4paper]{amsart}
\usepackage[margin=19mm]{geometry}
\usepackage{amsmath,amssymb,mathtools}
\usepackage[scr=rsfs,cal=euler]{mathalfa}
\usepackage{xfrac}
\usepackage[unicode,hidelinks,bookmarks=false]{hyperref}
\newcommand{\ie}{i.e.\ }
\newcommand{\cf}{cf.\ }
\newcommand{\resp}{respectively\ }
\newcommand{\Subsection}{\subsection}
\providecommand{\nobreakdash}{\nobreak\hbox{-}}
\setlength{\emergencystretch}{2em}
\multlinegap0pt
\usepackage{amssymb}
\usepackage{dsfont}
\usepackage{float}
\usepackage{mdframed}
\usepackage{mathtools}
\usepackage{xcolor}
\usepackage[shortlabels]{enumitem}
\newtheorem{proposition}{Proposition}
\renewcommand{\Re}{\operatorname{Re}\,}
\title{Dyson Brownian motion on a Jordan curve}
\author{Vladislav Guskov, Mingchang Liu, Fredrik Viklund}
\date{Source: arXiv:2603.05316v1 (CC BY 4.0)}
\begin{document}
\maketitle

\noindent\emph{Source and licence.} Dyson Brownian motion on a Jordan curve. Version arXiv:2603.05316v1 (CC BY 4.0), distributed by arXiv under the Creative Commons Attribution 4.0 International licence, http://creativecommons.org/licenses/by/4.0/, as recorded in the arXiv licence metadata for this exact version and shown on the abstract page https://arxiv.org/abs/2603.05316v1. Full licence notice: \texttt{licenses/arxiv-2603.05316-CC-BY-4.0.txt}.\par\smallskip

\noindent\textit{Editorial scope.} Counted unit: the article's proposition prop::Kolmogorov\_diffusion (source0.tex lines 940--943). The article's complete original proof of it is delivered with it (source0.tex lines 944--1045, one \texttt{proof} environment of 102 lines). No mathematics is written, repaired or re-derived: the statement and the proof are the article's own lines, in the article's own notation, and no retained line is substituted or re-flowed.  Retained uncounted as the source-local context the proof needs: lines 822--826; lines 859--874.  Nothing was omitted from any retained span, so the package carries no omission record.  No retained line contains a citation, so the wrapper carries no bibliography and no bibliography entry.  No retained line contains a figure environment, an image-inclusion command or drawing code, so no image is delivered and the package declares no asset dependency.  Editorial formatting only, in addition: an AMS \texttt{amsart} preamble that restates the article's own declarations and macros used by the retained lines, namely the environments \texttt{proposition} on separate counters, together with the standard packages those lines need.  The retained environments print in this document as Proposition 1 in order of appearance, whereas the article prints the same items with the numbering of its own full document.  The complete native source of the article as distributed in the arXiv e-print for this exact version is bundled unchanged as \texttt{sources/195792/source0.tex}.
\begin{equation}\label{eq::Cauchy_probem_x}
    \int_{\Omega}\varphi d\mu_{t} 
    = \varphi(x) + \int_{0}^{t}\int_{\Omega}L\varphi d\mu_{s}ds
    \quad\forall \varphi\in C^{\infty}_{0}(\Omega).
\end{equation}    

\begin{proposition}\label{prop::FPK_X_process}
The transition probability function $P(x,t,dy)$ of the parametrization process $X^{x}$ solves the Cauchy problem \eqref{eq::Cauchy_probem_x} with $\rho=\rho_{\beta, N}$. Moreover, the transition probability measure is absolutely continuous with respect to the Lebesgue measure 
\begin{align*}
    P(x,t,dy) = p(x,t,y)dy,
\end{align*}
$(t,y)\to p(x,t,y)$ is positive and locally Hölder continuous on $(0,+\infty)\times D$, and the density satisfies the Fokker-Planck-Kolmogorov equation with reflecting boundary conditions
\begin{align*}
    &\partial_{t}p(x,t,\bullet)=L^{*}p(x,t,\bullet) \text{ weakly in }D,\\
    &n\cdot\left(-\frac{1}{2}\rho_{\beta, N}\nabla \frac{p}{\rho_{\beta, N}} \right) = 0 \text{ weakly on }\partial D.
\end{align*}
% In addition, the ratio
% \[
%    (t,x)\mapsto \sup_{y\in D}\frac{p(x,t,y)}{\rho_{\beta, N}(y)}
% \]
% is locally bounded on $(0,+\infty)\times D$.
\end{proposition}

\begin{proposition}\label{prop::Kolmogorov_diffusion}
Suppose $\gamma \in C^{2}$.
Then, the parametrization process $X^{x}$ is a diffusion process in the sense of Kolmogorov with identity diffusion matrix and the drift coefficient given by $b =\frac{1}{2}\nabla\log\rho_{\beta, N}$.   
\end{proposition}

\begin{proof}
$1^{\circ}.$ Let $f\in C^{\infty}(D)$. The time derivative of $(P_{t}f)(x)=\mathbb{E}[X^{x}(t)]$ at $t=0$ can be written as 
    \begin{align*}
        \frac{d}{dt}\int_{ D}f(y)P(x,t,dy)\Big|_{t=0}
        &= \lim\limits_{h\to 0+}\frac{1}{h}\left(\int_{ D}f(y)P(x,h,dy) - \int_{ D}f(y')P(x,0,dy')\right) \\ 
        &=\lim\limits_{h\to 0+}\frac{1}{h}\left(\int_{ D}f(y)P(x,h,dy) -f(x)\right)  \\
        & = \lim\limits_{h\to 0+}\frac{1}{h}\int_{ D}\left(f(y)-f(x)\right)P(x,h,dy) .
    \end{align*}
On the other hand, as in the proof of Proposition~\ref{prop::FPK_X_process}, from Itô's formula we have
\begin{align*}
     \frac{d}{dt}\int_{ D}f(y)P(x,t,dy)\Big|_{t=0}
     = \int_{D}(Lf)(y)P(x,t,dy)\Big|_{t=0}
     = (Lf)(x).
\end{align*}
We obtain the identity
\[
    \lim\limits_{h\to 0+}\frac{1}{h}\int_{ D}\left(f(y)-f(x)\right)P(x,h,dy)  = (Lf)(x) \quad\forall f\in C^{\infty}(D).
\]
Taking $f_{k}(y)=y_{k}$ yields
\begin{align*}
\lim\limits_{h\to 0+}\frac{1}{h}\int_{ D}\left(y_{k}-x_{k}\right)P(x,h,dy) 
=(Lf_{k})(x)=b_{k}(x).
\end{align*}
Taking $f_{ij}(y) =(y_{i}-x_{i})(y_{j}-x_{j})z_{i}z_{j}$ yields
\begin{align*}
   \lim\limits_{h\to 0+}\frac{1}{h}\int_{ D}(y_{i}-x_{i})(y_{j}-x_{j})z_{i}z_{j}P(x,h,dy) =  (Lf_{ij})(x) = z_{i}^2\delta_{ij}.
\end{align*}

$2^{\circ}.$ Next we show stochastic continuity of the transition probability function, that is, for all $x\in D$ and $\delta > 0$ such that $\overline{B(x,\delta)}\subset D$ the following limit holds
\begin{align*}
    \lim\limits_{h\to 0+}\frac{1}{h}P(x,h,B(x,\delta)^{c})=0.
\end{align*}
By Itô's formula, for $p>2$,
\begin{align*}
    \mathbb{E}\left[|X^{x}(t)-x|^{p}\right] 
    = &p\frac{N+p-2}{2}\int_{0}^{t}\mathbb{E}\left[|X^{x}(s)-x|^{p-2}\right]ds \\
    &+ p \int_{0}^{t}\mathbb{E}\left[|X^{x}(s)-x|^{p-2}\sum\limits_{i=1}^{N}b_{i}(X^{x}(s))((X^{x}(s))_{i}-x_{i})\right]ds.
\end{align*}
The sum in the second term can be written as
\begin{align*}
    \sum\limits_{i=1}^{N}b_{i}(y)(y_{i}-x_{i}) 
    = \sum\limits_{i=1}^{N}(b_{i}(y)-b_{i}(x))(y_{i}-x_{i})  +\sum\limits_{i=1}^{N} (y_{i}-x_{i})b_{i}(x). 
\end{align*}
Now we will show that the first term on the right hand side is bounded from above.
\begin{align*}
     &\sum\limits_{i=1}^{N}(b_{i}(y)-b_{i}(x))(y_{i}-x_{i}) \\
     &= \frac{\beta}{2}\sum\limits_{i=1}^{N}(y_{i}-x_{i})\sum\limits_{j:j\neq i}\Re\left\{\frac{\gamma'(y_{i})}{\gamma(y_{i})-\gamma(y_{j})}-\frac{\gamma'(x_{i})}{\gamma(x_{i})-\gamma(x_{j})} \right\}\\
     &= \frac{\beta}{2}\sum_{i < j}\Re\left\{\frac{(y_{i}-x_{i})\gamma'(y_{i})-(y_{j}-x_{j})\gamma'(y_{j})}{\gamma(y_{i})-\gamma(y_{j})}\right\}
    + \frac{\beta}{2} \sum_{i< j}\Re\left\{\frac{(x_{i}-y_{i})\gamma'(x_{i})-(x_{j}-y_{j})\gamma'(x_{j})}{\gamma(x_{i})-\gamma(x_{j})}\right\}.
\end{align*}
Now we take a closer look at the first term, which can be factored as
\begin{align}
    &\Re\left\{\frac{(y_{i}-x_{i})\gamma'(y_{i})-(y_{j}-x_{j})\gamma'(y_{j})}{\gamma(y_{i})-\gamma(y_{j})}\right\} \\
    &= (y_{i}-x_{i})\Re\left\{\frac{\gamma'(y_{i})-\gamma'(y_{j})}{\gamma(y_{i})-\gamma(y_{j})}\right\}
    + \left(1 - \frac{x_{j}-x_{i}}{y_{j}-y_{i}}\right)\Re\left\{\gamma'(y_{j})\frac{y_{i}-y_{j}}{\gamma(y_{i})-\gamma(y_{j})}\right\}.\label{eq::monotone_drift_diverges}
\end{align}
Since $x,y\in D$, it holds that $x_{i}<x_{j}$ and $y_{i}<y_{j}$ for $i<j$, and the term $\frac{x_{i}-x_{j}}{y_{i}-y_{j}}$ is always positive.  If $y_{j}-y_{i}\to0+$, the expression \eqref{eq::monotone_drift_diverges} can only go to $-\infty$, because under the assumption $\gamma\in C^2$
\[
    \lim\limits_{y_{j}-y_{i}\to 0+}\Re\left\{\gamma'(y_{j})\frac{y_{i}-y_{j}}{\gamma(y_{i})-\gamma(y_{j})}\right\} = 1
\]
and 
\[
     \lim\limits_{y_{j}-y_{i}\to 0+}  \Re\left\{\frac{\gamma'(y_{i})-\gamma'(y_{j})}{\gamma(y_{i})-\gamma(y_{j})}\right\} 
     = \Re\left\{\frac{\gamma''(y_{j})}{\gamma'(y_{j})}\right\} 
     =  \Re\left\{\gamma''(y_{j}) \overline{\gamma'(y_{j})}\right\} 
     =0.
\]
The same conclusion holds when $y_{j}-y_{i}\to l$ since the expression \eqref{eq::monotone_drift_diverges} can be rewritten as
\begin{align*}
    &\Re\left\{\frac{(y_{i}-x_{i})\gamma'(y_{i})-(y_{j}-x_{j})\gamma'(y_{j})}{\gamma(y_{i})-\gamma(y_{j})}\right\} \\
    &= (y_{i}-x_{i})\Re\left\{\frac{\gamma'(y_{i})-\gamma'(y_{j})}{\gamma(y_{i})-\gamma(y_{j})}\right\}+ \left(1 - \frac{l+x_{i}-x_{j}}{l+y_{i}-y_{j}}\right)\Re\left\{\gamma'(y_{j})\frac{l+y_{i}-y_{j}}{\gamma(y_{i})-\gamma(y_{j})}\right\}.
\end{align*}
Hence, 
\begin{align*}
    \lim\limits_{y\to \partial D}(b(x)-b(y))\cdot (x-y)= -\infty,
\end{align*}
and similarly for $x\to\partial D$. There exist a constant $C=C(\gamma, N)>0$ such that 
\begin{align*}
    \sum\limits_{i=1}^{N}(b_{i}(y)-b_{i}(x))(y_{i}-x_{i}) < C(\gamma, N).
\end{align*}
Hence, there are positive constants $c_{1},c_{2}$, that depend on $N,p, C(\gamma, N), |b(x)|, \beta $, such that 
\begin{align*}
    \mathbb{E}\left[|X^{x}(h)-x|^{p}\right] 
    \le c_{1} \int_{0}^{h}\mathbb{E}\left[|X^{x}(s)-x|^{p-2}\right]ds 
    + c_{2} \int_{0}^{h}\mathbb{E}\left[|X^{x}(s)-x|^{p-1}\right]ds.
\end{align*}
This implies that 
\[
    \lim\limits_{h\to 0+}\frac{1}{h}\mathbb{E}\left[|X^{x}(h)-x|^{p}\right] =0.
\]
By Chebyshev's inequality
\[
    P(x,h,B(x,\delta)^{c}) 
    = \mathbb{P}[|X^{x}(h)-x|\ge \delta]
    \le \delta^{-p}\mathbb{E}\left[|X^{x}(h)-x|^{p}\right].
\]
Hence, for all $x\in D$ and $\delta > 0$ such that $\overline{B(x,\delta)}\subset D$ we have
\begin{align*}
    \lim\limits_{h\to 0+}\frac{1}{h}P(x,h,B(x,\delta)^{c})=0.
\end{align*}
Consequently, combining the stochastic continuity with the results of $1^{\circ}.$ gives items $1.$ and $2.$ of the proposition. The fact that the equations $1.-3.$ can be written in terms of densities of the transition probabilities $P(x,t,dy)$ follows from Proposition~\ref{prop::FPK_X_process}.
\end{proof}

\end{document}
